\documentclass[11pt,a4paper]{article}
\usepackage[T1]{fontenc}
\usepackage[utf8]{inputenc}
\usepackage{lmodern}
\usepackage[margin=2.35cm]{geometry}
\usepackage{amsmath,amssymb,amsthm}
\usepackage{microtype}
\usepackage[hidelinks]{hyperref}

\hypersetup{
  pdftitle={A Square-Root-Two Exponent for Bilinear Sum-of-Squares Composition Formulas},
  pdfsubject={A conditional algebraic consequence of the crossing decomposition in Appendix B of a 3-SAT manuscript},
  pdfkeywords={sum of squares, Hurwitz problem, crossing decomposition, bilinear forms}
}
\newtheorem{theorem}{Theorem}
\newcommand{\CC}{\mathbb C}
\newcommand{\QQ}{\mathbb Q}
\newcommand{\RR}{\mathbb R}
\newcommand{\ii}{\mathrm i}
\setlength{\parindent}{0pt}
\setlength{\parskip}{0.5em}
\setlength{\emergencystretch}{1.2em}
\allowdisplaybreaks[1]

\title{\textbf{An \(N^{\sqrt{2}+o(1)}\) Upper Bound for Bilinear\\
Sum-of-Squares Composition Formulas}\\[0.7em]
\large An algebraic consequence of Appendix B of\\
\emph{A \(2^{o(n)}\) Time Algorithm for 3-SAT}}
\author{}
\date{27 October 2026}

\begin{document}
\maketitle

\begin{abstract}
Assuming the algebraic identities and uniform rank estimates in Appendix~B
of the 26 October 2026 Claude/Anthropic manuscript \emph{A \(2^{o(n)}\)
Time Algorithm for 3-SAT}, its crossing decomposition implies an upper
bound of \(N^{\sqrt{2}+o(1)}\) for Naslund's crossing-decomposition quantity
\(P(N)\). It also yields an identity expressing the product of two sums of
\(N\) squares as a sum of at most \(N^{\sqrt{2}+o(1)}\) products of rational
bilinear forms. Over the complex numbers, each such product is a sum of
two squares of bilinear forms. Consequently, the same hypotheses imply a
complex bilinear sum-of-squares composition formula with at most
\(N^{\sqrt{2}+o(1)}\) squares, improving the \(O(N^{1.62})\) bound of
Hrube\v{s} (2024, 2026). No modification of the construction in Appendix~B
is necessary. The implication uses only its algebraic identities and rank
estimates, not the SAT algorithm or its running-time analysis.
\end{abstract}

\section{Main result}

Let \(M_{\CC}(N)\) denote the minimum integer \(M\) for which there exist
complex bilinear forms \(F_1,\ldots,F_M\) satisfying the polynomial identity
\[
  \left(\sum_{i=1}^{N}a_i^2\right)
  \left(\sum_{j=1}^{N}b_j^2\right)
  =\sum_{\ell=1}^{M}F_\ell(a,b)^2.
\]
Here each form is homogeneous of degree one in \(a\) and of degree one in
\(b\). The squares are algebraic squares, without complex conjugation.

\begin{theorem}[Conditional on the Appendix B construction]
\label{thm:main}
Assume the identities and the uniform rank-factorization estimate
\eqref{eq:decomposition}--\eqref{eq:rank} below, established in
Appendix~B of \emph{A \(2^{o(n)}\) Time Algorithm for 3-SAT}
\cite{claude}. There is an absolute constant \(C\) such that, for every
\(N\geq 2\),
\[
  M_{\CC}(N)=O\!\left(N^{\sqrt{2}}(\log N)^C\right).
\]
In particular, \(M_{\CC}(N)\leq N^{\sqrt{2}+o(1)}\).
The same upper bound holds for the number of rational bilinear products
required to express the product of two sums of squares, and for the
crossing-decomposition quantity \(P(N)\) defined in Appendix~G of that
manuscript.
\end{theorem}

\section{The crossing decomposition}

First suppose \(N=2^\lambda\), and identify \([N]=\{1,\ldots,N\}\) with
\(\{0,1\}^{\lambda}\). Define the crossing tensor
\[
  D(x,y,u,v)=[x=v][y=u],
\]
where brackets denote indicators.

Appendix~B, Theorem~B, constructs an integer \(h=O(\lambda^2)\),
integer-valued functions
\[
  A_s(x,y,u,v),\qquad -2h\leq s\leq 2h,
\]
and a nonzero integer \(d=h!\) such that
\begin{equation}
  dD(x,y,u,v)=\sum_{s=-2h}^{2h}A_s(x,y,u,v).
  \label{eq:decomposition}
\end{equation}
Moreover, each \(A_s\) admits two factorizations. The first separates
\((x,y)\) from \((u,v)\):
\begin{equation}
  A_s(x,y,u,v)
  =\sum_{q=1}^{R_s}f_{s,q}(x,y)g_{s,q}(u,v).
  \label{eq:row}
\end{equation}
The second separates \((x,u)\) from \((y,v)\):
\begin{equation}
  A_s(x,y,u,v)
  =\sum_{q=1}^{C_s}
    \widetilde f_{s,q}(x,u)\widetilde g_{s,q}(y,v).
  \label{eq:column}
\end{equation}
The factors are integer-valued, and the positive factorization lengths
satisfy
\begin{equation}
  \log_2 R_s+\log_2 C_s
  \leq 2\sqrt{2}\,\lambda+O\!\left(\log(2\lambda)\right),
  \label{eq:rank}
\end{equation}
uniformly in \(s\). This follows from the bound on the rounded-up binary
logarithms in Appendix~B, equation~(56), derived using the Hardy--Ramanujan
asymptotics for integer partitions. The \(O(\log(2\lambda))\) term agrees
with \(O(\log\lambda)\) asymptotically and also covers \(\lambda=1\).

\section{Choosing the cheaper factorization}

For each \(s\), use whichever of \eqref{eq:row} and \eqref{eq:column}
contains fewer terms. By \eqref{eq:rank},
\[
  \min(R_s,C_s)
  \leq \sqrt{R_sC_s}
  \leq 2^{\sqrt{2}\lambda+O(\log(2\lambda))}.
\]
Since \(N=2^\lambda\), this gives
\begin{equation}
  \min(R_s,C_s)
  \leq N^{\sqrt{2}}(\log N)^{O(1)}.
  \label{eq:minimum}
\end{equation}

There are \(4h+1=O(\lambda^2)\) values of \(s\). Expanding each \(A_s\)
using its cheaper orientation and dividing \eqref{eq:decomposition} by
\(d\) therefore expresses \(D\) as a sum of \(L\) rational rank-one
functions, each separated in one of the two orientations, where
\begin{equation}
  L=O\!\left(N^{\sqrt{2}}(\log N)^C\right)
  \label{eq:length}
\end{equation}
for an absolute constant \(C\). By the definition in Appendix~G, this proves
\begin{equation}
  P(N)\leq L\leq N^{\sqrt{2}+o(1)}.
  \label{eq:crossing-bound}
\end{equation}

\section{From crossing factorizations to bilinear products}

Multiply \eqref{eq:decomposition}, after division by \(d\), by
\(a_xb_ya_vb_u\), and sum over all \(x,y,u,v\in[N]\). The left-hand side
becomes
\begin{equation}
  \sum_{x,y,u,v}[x=v][y=u]\,a_xb_ya_vb_u
  =\left(\sum_xa_x^2\right)\left(\sum_yb_y^2\right).
  \label{eq:contraction}
\end{equation}

For a rank-one term in orientation \eqref{eq:row}, namely
\(f(x,y)g(u,v)\), the corresponding contribution, before normalization by
\(d\), is
\[
  \left(\sum_{x,y}f(x,y)a_xb_y\right)
  \left(\sum_{u,v}g(u,v)b_ua_v\right),
\]
a product of two bilinear forms. For a rank-one term in orientation
\eqref{eq:column}, namely
\(\widetilde f(x,u)\widetilde g(y,v)\), the contribution is
\[
  \left(\sum_{x,u}\widetilde f(x,u)a_xb_u\right)
  \left(\sum_{y,v}\widetilde g(y,v)b_ya_v\right).
\]
This is again a product of two bilinear forms.

Hence \eqref{eq:decomposition} gives an identity
\begin{equation}
  \boxed{
  \left(\sum_{i=1}^{N}a_i^2\right)
  \left(\sum_{j=1}^{N}b_j^2\right)
  =\sum_{q=1}^{L}F_q(a,b)G_q(a,b),
  }
  \label{eq:products}
\end{equation}
where all \(F_q,G_q\) are rational bilinear forms and
\(L\leq N^{\sqrt{2}+o(1)}\). The normalization factor \(1/h!\) is absorbed
into one factor of each product.

\section{Conversion to complex sums of squares}

For any bilinear forms \(F,G\) over \(\QQ\), the identity
\begin{equation}
  FG=
  \left(\frac{F+G}{2}\right)^2
  +\left(\frac{\ii(F-G)}{2}\right)^2
  \label{eq:polarization}
\end{equation}
holds over \(\CC\), where \(\ii^2=-1\). Both expressions inside the
squares remain bilinear forms in \(a,b\).

Applying \eqref{eq:polarization} to every term of \eqref{eq:products}
produces a sum of at most \(2L\) squares of complex bilinear forms. Thus
\[
  M_{\CC}(N)\leq 2L
  =O\!\left(N^{\sqrt{2}}(\log N)^C\right).
\]
In fact, the coefficients produced lie in \(\QQ(\ii)\).

For arbitrary \(N\), choose \(\lambda=\lceil\log_2N\rceil\), apply the
construction to \(2^\lambda\) variables, and set the excess variables to
zero. Since \(N\leq 2^\lambda<2N\), the same asymptotic bound follows.
Restricting all four tensor indices to \([N]\) likewise gives the bound
for \(P(N)\). This proves Theorem~\ref{thm:main}.

\section{Comparison with prior work}

Hrube\v{s} established an \(O(N^{1.62})\) upper bound for complex bilinear
sum-of-squares composition formulas in \emph{A Subquadratic Upper Bound on
Sum-Of-Squares Composition Formulas} (CCC 2024) \cite{hrubes2024},
subsequently developed in \emph{A Subquadratic Upper Bound on Hurwitz's
Problem and Related Noncommutative Polynomials} (2026)
\cite{hrubes2026}. Related connections to noncommutative circuits appear
in \cite{hwy2011}.

The construction above gives
\[
  \sqrt{2}\approx 1.41421356<1.62,
\]
so, under the stated hypotheses, it improves the exponent by
approximately \(0.206\).

This conclusion is conditional only on the correctness of the algebraic
identities and uniform rank-factorization estimates in the
Claude/Anthropic manuscript. It does not rely on that manuscript's SAT
algorithm, recursive running-time analysis, or formalization of the
Exponential Time Hypothesis (ETH). This note proves the implication from
the stated algebraic hypotheses; it does not independently establish
those hypotheses or the claimed SAT running time.

\section{Remarks and attribution}

\paragraph{Rational versus complex coefficients.}
Equation~\eqref{eq:products} holds over \(\QQ\), as the Appendix~B factors
are integer-valued before normalization.
Equation~\eqref{eq:polarization} uses \(\ii=\sqrt{-1}\), so the asserted
upper bound for sums of squares is over \(\CC\). It does not establish a
sum-of-squares identity of the same size over \(\RR\) or \(\QQ\).

\paragraph{Relation to Naslund's problem.}
Equation~\eqref{eq:crossing-bound} strengthens, under the stated
hypotheses, the upper bound for \(P(N)\) discussed in Appendix~G.
Naslund's partition-rank work is \cite{naslund2020}. The conversion to
bilinear products and then complex squares establishes a separate
consequence for Hurwitz's composition problem.

\paragraph{Attribution.}
The construction is that of Appendix~B of Claude's manuscript for
Anthropic. The observation here is that its two-sided
factorization estimate implies a stronger bound for \(P(N)\) and for
complex sum-of-squares composition formulas.

\begin{thebibliography}{9}
\setlength{\itemsep}{0.35em}

\bibitem{claude}
Claude (Anthropic).
\emph{A \(2^{o(n)}\) Time Algorithm for 3-SAT}.
Manuscript dated 26 October 2026.
Appendix~B, Theorem~B, equations (53)--(56); Appendix~G.
\href{https://github.com/anthropics/formal-math/blob/16a5ce41304fcdaea4e0be7c132e1c361dc26da5/eth/paper.pdf}
{Version in \texttt{anthropics/formal-math}, commit \texttt{16a5ce4}}.

\bibitem{hrubes2024}
Pavel Hrube\v{s}.
\emph{A Subquadratic Upper Bound on Sum-Of-Squares Composition Formulas}.
In \emph{39th Computational Complexity Conference (CCC 2024)},
LIPIcs 300, Article 12, pp.~12:1--12:11, 2024.
\href{https://doi.org/10.4230/LIPIcs.CCC.2024.12}
{doi:10.4230/LIPIcs.CCC.2024.12}.

\bibitem{hrubes2026}
Pavel Hrube\v{s}.
\emph{A Subquadratic Upper Bound on Hurwitz's Problem and Related
Noncommutative Polynomials}.
\emph{SIAM Journal on Computing} 55(2), 451--467, 2026.
\href{https://doi.org/10.1137/24M1666884}{doi:10.1137/24M1666884}.

\bibitem{hwy2011}
Pavel Hrube\v{s}, Avi Wigderson, and Amir Yehudayoff.
\emph{Non-commutative Circuits and the Sum of Squares Problem}.
\emph{Journal of the American Mathematical Society} 24(3), 871--898, 2011.
\href{https://doi.org/10.1090/S0894-0347-2011-00694-2}
{doi:10.1090/S0894-0347-2011-00694-2}.

\bibitem{naslund2020}
Eric Naslund.
\emph{The Partition Rank of a Tensor and \(k\)-Right Corners in
\(\mathbb F_q^n\)}.
\emph{Journal of Combinatorial Theory, Series A} 174, Article 105190, 2020.
\href{https://doi.org/10.1016/j.jcta.2019.105190}
{doi:10.1016/j.jcta.2019.105190}.

\end{thebibliography}
\end{document}
